% !TeX root = main.tex
% =============================================================================
% SEÇÃO 3: ATIVIDADES EM ANDAMENTO E FUTURAS
% =============================================================================
% Esta seção deve apresentar:
% - Atividades que estão sendo desenvolvidas no momento atual
% - Planejamento das atividades futuras até o final do projeto
% - Cronograma previsto (se aplicável)
% - Objetivos esperados para cada atividade
% - Recursos e metodologias que serão utilizados
% =============================================================================

\section{Atividades em Andamento / Futuras}

% Introdução sobre o planejamento das próximas etapas
A seguir, estão sintetizadas as principais atividades que deverão ser 
realizadas durante o período de [mês/ano] a [mês/ano]:

\begin{enumerate}
    % --- Atividade em Andamento/Futura 1 ---
    \item [\textbf{i.}] \textbf{[Título da Primeira Atividade Planejada]:} 
    [Descreva a primeira atividade planejada ou em andamento, incluindo 
    seus objetivos, metodologia prevista, recursos necessários e resultados 
    esperados. Se a atividade depende de aprovações ou etapas anteriores, 
    mencione essas dependências.]
    
    % --- Atividade em Andamento/Futura 2 ---
    \item [\textbf{ii.}] \textbf{[Título da Segunda Atividade Planejada]:} 
    [Descreva a segunda atividade, destacando como ela se relaciona com 
    as atividades anteriores e quais são os resultados esperados. Seja 
    específico quanto aos métodos e abordagens que serão utilizados.]
    
    % Exemplo: A análise dos dados será realizada seguindo a metodologia
    % proposta por \citeonline{autor2023_exemplo}, utilizando as ferramentas
    % recomendadas pela literatura \cite{autor2022_livro,organizacao2024_manual}.
    
    % --- Atividade em Andamento/Futura 3 ---
    \item [\textbf{iii.}] \textbf{[Título da Terceira Atividade Planejada]:} 
    [Continue descrevendo as atividades futuras. Se esta atividade envolve 
    a produção de resultados científicos (artigos, relatórios, apresentações), 
    mencione os veículos de divulgação previstos e os prazos esperados.]
    
\end{enumerate}

% Nota: Você pode adicionar uma subseção para apresentar um cronograma
% mais detalhado, se necessário.

% --- Cronograma Resumido (Opcional) ---
% \subsection{Cronograma Previsto}
% 
% [Apresente aqui um cronograma detalhado com as datas previstas para 
% cada atividade, podendo utilizar uma tabela ou lista estruturada.]